% ======================================================================
% capitulos/ce02_ingenieria_software.tex
% ----------------------------------------------------------------------
% LÍNEA CE02 — INGENIERÍA DE SOFTWARE
% Esqueleto / marcador de posición. Reemplazar por los capítulos reales
% conforme se desarrollen los entregables de esta línea.
% ======================================================================
\chapter{Ingeniería de Software — Presentación de la línea}
\label{chap:ce02-intro}

Esta línea (CE02) desarrolla la especificación, el diseño y la construcción del
software del \textbf{Sistema de Gestión Integral para ISP de Fibra Óptica
Multi-Zona}. Toma como insumo los entregables de la línea CE01 (en especial el
\hyperref[chap:procesos]{modelado de procesos (TO-BE)} y la
\hyperref[chap:soluciontic]{solución TIC}) y produce los artefactos de ingeniería
que hacen construible el sistema.

\begin{upeubox}{Entregables planificados de la línea CE02 (en desarrollo)}
Según la EDT del \hyperref[chap:plangestion]{Plan de Gestión}, esta línea comprende:
\begin{itemize}[leftmargin=1.4em]
  \item \textbf{Especificación de requisitos (SRS):} requerimientos funcionales y no
  funcionales conforme a IEEE~29148.
  \item \textbf{Prototipos navegables:} diseño de interfaz (Figma) validado con WISP.
  \item \textbf{Arquitectura de software:} vistas C4 e IEEE/IEEE~42010, y modelado
  UML (casos de uso, clases, secuencia, despliegue).
  \item \textbf{Diseño de base de datos:} modelo entidad-relación, modelo lógico,
  diccionario de datos y scripts SQL (DDL/DML).
  \item \textbf{Construcción del sistema (semestre 2):} backend Django (API REST),
  frontend web Angular, app móvil e integración MikroTik.
\end{itemize}
\end{upeubox}

\medskip
\textit{Nota: los capítulos de esta línea se incorporarán aquí a medida que se
completen sus entregables. Este archivo es un marcador de posición para mantener
integrada la estructura de las tres líneas del proyecto.}

% Al desarrollar, sustituir este archivo (o añadir \input adicionales en main.tex)
% por capítulos como:
%   \input{capitulos/ce02_srs}
%   \input{capitulos/ce02_prototipos}
%   \input{capitulos/ce02_arquitectura_uml}
%   \input{capitulos/ce02_base_datos}
%   \input{capitulos/ce02_construccion}
